\documentclass[10pt,a4paper]{amsart}
\usepackage[margin=19mm]{geometry}
\usepackage{amsmath,amssymb,mathtools}
\usepackage[scr=rsfs,cal=euler]{mathalfa}
\usepackage{xfrac}
\usepackage[unicode,hidelinks,bookmarks=false]{hyperref}
\newcommand{\ie}{i.e.\ }
\newcommand{\cf}{cf.\ }
\newcommand{\resp}{respectively\ }
\newcommand{\Subsection}{\subsection}
\providecommand{\nobreakdash}{\nobreak\hbox{-}}
\setlength{\emergencystretch}{2em}
\multlinegap0pt
\usepackage{tikz}
\usepackage{amssymb}
\usepackage{enumerate}
\usepackage{mathtools}
\newtheorem{theorem}{Theorem}
\newtheorem{example}{Example}
\newtheorem{proposition}{Proposition}
\DeclareMathOperator{\Aut}{{\rm Aut}}
\DeclareMathOperator{\DU}{\Delta {\rm U}}
\newcommand{\FF}{\mathbb{F}}
\newcommand{\Fl}{\mathbb{F}_\ell}
\DeclareMathOperator{\GL}{{\rm GL}}
\newcommand{\GQr}{G_{\QQ(\zeta_{r})}}
\newcommand{\GQrl}{G_{\QQ(\zeta_{r\ell})}}
\DeclareMathOperator{\GU}{{\rm GU}}
\newcommand{\Glam}{G_{\lambda}}
\newcommand{\QQ}{\mathbb{Q}}
\newcommand{\Qr}{{\QQ(\zeta_r)}}
\DeclareMathOperator{\SU}{{\rm SU}}
\newcommand{\p}{\mathfrak{p}}
\title{Superelliptic curves with large Galois images}
\author{Pip Goodman}
\date{Source: arXiv:2011.14461v2 (CC BY 4.0)}
\begin{document}
\maketitle

\noindent\emph{Source and licence.} Superelliptic curves with large Galois images. Version arXiv:2011.14461v2 (CC BY 4.0), distributed by arXiv under the Creative Commons Attribution 4.0 International licence, http://creativecommons.org/licenses/by/4.0/, as recorded in the arXiv licence metadata for this exact version and shown on the abstract page https://arxiv.org/abs/2011.14461v2. Full licence notice: \texttt{licenses/arxiv-2011.14461-CC-BY-4.0.txt}.\par\smallskip

\noindent\textit{Editorial scope.} Counted unit: the article's proposition prop\_DU\_image (source0.tex lines 1960--1967). The article's complete original proof of it is delivered with it (source0.tex lines 1970--1996, one \texttt{proof} environment of 27 lines). No mathematics is written, repaired or re-derived: the statement and the proof are the article's own lines, in the article's own notation, and no retained line is substituted or re-flowed.  Retained uncounted as the source-local context the proof needs: lines 554--559; lines 1805--1813.  Nothing was omitted from any retained span, so the package carries no omission record.  No retained line contains a citation, so the wrapper carries no bibliography and no bibliography entry.  No retained line contains a figure environment, an image-inclusion command or drawing code, so no image is delivered and the package declares no asset dependency.  Editorial formatting only, in addition: an AMS \texttt{amsart} preamble that restates the article's own declarations and macros used by the retained lines, namely the environments \texttt{theorem}, \texttt{example}, \texttt{proposition} on separate counters, together with the standard packages those lines need.  The retained environments print in this document as Theorem 1, Example 1, Proposition 1 in order of appearance, whereas the article prints the same items with the numbering of its own full document.  The complete native source of the article as distributed in the arXiv e-print for this exact version is bundled unchanged as \texttt{sources/105222/source0.tex}.
\begin{theorem}
\label{thm_containment_of_the_image_of_GQr_and_GQrl}
The image of $\rho_\lambda \colon \GQr \rightarrow \Aut(J[\lambda])$ is contained in $\GL_n(\ell^i)$ if $i$ is odd and $\DU_n(\ell^{i/2})$ if $i$ is even.

Moreover, the subgroup $G_\lambda = \rho_\lambda(\GQrl)$ is contained in $\GL(\ell^i)$ if $i$ is odd and $\GU_n(\ell^{i/2})$ if $i$ is even.
\end{theorem}

\begin{example}
\label{inertia_at_l_for_primes_of_res_deg_1_and_2}
Let $\ell \equiv 1 \mod{r}$ be a prime of semistable reduction for $J$. Then $\ell$ is totally split in $\Qr$. Let $\lambda,\lambda'$ be primes above $\ell$ in $\Qr$ and let $1\leq j \leq r-1$ be such that $\psi_j(\lambda') = \lambda$. Then, as the unique fundamental character of level one is the mod $\ell$ cyclotomic character $\chi_\ell$, we have
$$\Bar{\Omega}_\lambda|_{I_{\lambda'}} = \chi_\ell^{\left\lfloor \frac{(r-j)d}{r}\right\rfloor_<} .$$

Let $\ell \equiv -1 \mod{r}$ be a prime of good reduction. We have $[k_\lambda:\Fl]=2$ for a prime $\lambda$ above $\ell$ in $\Qr$. Let $\lambda'$ be a prime above $\ell$ in $\Qr$ and let $1 \leq j \leq r-1$ be such that $\psi_j(\lambda') = \lambda$. Then
$$\Bar{\Omega}_\lambda|_{I_{\lambda'}} = \theta^{\left\lfloor \frac{(r-j)d}{r}\right\rfloor_< + \ell\left\lfloor \frac{jd}{r}\right\rfloor_<} $$
where $\theta$ is some fundamental character of level 2.
\end{example}

\begin{proposition}
\label{prop_DU_image}
Suppose $2r|d$. Let $\ell \equiv -1 \mod{r}$, $\ell \neq 2$ be a prime of semistable reduction for $J/\Qr$.

Suppose there exists some $\delta|2r$ such that $\gcd(\frac{d}{r},\frac{\ell+1}{\delta})= \gcd(\delta, \frac{\ell+1}{\delta})=1$ and $f$ has $\p$-degree 2 for some prime of $\Qr$ with residue characteristic $p \neq r,\ell$.

If $\rho_\lambda(\GQr)$ contains $\SU_n(\ell)$, then $\rho_\lambda(\GQr) = \DU_n(\ell)$.
\end{proposition}

\begin{proof}
By Theorem \ref{thm_containment_of_the_image_of_GQr_and_GQrl} $\rho_\lambda(\GQr)$ is contained in a group isomorphic to $ \DU_n(\ell)$.
Thus it suffices to show $\det \circ \rho_\lambda = \Bar{\Omega}_\lambda : \GQr \rightarrow \FF_{\ell^2}^*$ is surjective. As $N^{\ell^2}_{\ell} \circ \Bar{\Omega}_\lambda$ coincides with the cyclotomic character, it suffices to show $\Bar{\Omega}_\lambda(\GQrl) = (\FF_{\ell^2}^*)^{\ell -1}$.

 Let $\delta|2r$ be a positive integer satisfying $\gcd(\delta, \frac{\ell+1}{\delta})=1$.
 As $f$ has $\p$-degree $2$, $\Bar{\Omega}_\lambda(I_\p) \leq (\FF^*_{\ell^2})^{\ell-1}$ has order $2r$.
It thus suffices to show $\Bar{\Omega}_\lambda(\GQrl)$ contains an element of order $\frac{\ell+1}{\delta}$.

Let $\lambda,\lambda'$ be primes above $\ell$ in $\Qr$ be such that $\psi_j(\lambda')=\lambda$. Example \ref{inertia_at_l_for_primes_of_res_deg_1_and_2} shows
\[
\Omega_\lambda|_{I_{\lambda'}} = \theta^{m_{j} + \ell m_{r-j} }
\]
for an appropriate choice of a level 2 fundamental character $\theta$ (this choice does not concern us as they differ by an automorphism of $\FF_{\ell^2}$).

As $\theta$ surjects onto $\FF_{\ell^2}^*$  and hence onto the kernel of the norm map, it suffices to compute  $\gcd(m_{j} + \ell m_{r-j} , \ell + 1)$ to determine the image of $I_{\lambda'}$ under $\Bar{\Omega}_\lambda$ which lands in $\Glam$.


Due to the congruence $m_{j} + \ell m_{r-j} \equiv m_{j} - m_{r-j}  \mod{ \ell + 1}$, we compute
\[m_{j} - m_{r-j} = \left\lfloor \frac{(r-j)d}{r}\right\rfloor_< - \left\lfloor \frac{jd}{r}\right\rfloor_< = \frac{d(r-2j)}{r}.
\]
This shows  $\frac{d}{r}$ divides $ m_{j} - m_{r-j}$ for all values of $j$. In fact, for $j = \frac{r-1}{2}$, we have $ m_{j} - m_{r-j} = \frac{d}{r}$, so this is the best contribution we will get from an inertia group at a prime above $\ell$.

The above shows $\Bar{\Omega}_\lambda(\GQrl)$ contains $(\FF^*_{\ell^2})^{\frac{d}{r}(\ell-1)}$.
Now let $\tau \in \FF_{\ell^2}^{\ell-1}$ be an element of order $\frac{\ell+1}{\delta}$.
By assumption $\gcd(\frac{d}{r},\frac{\ell+1}{\delta})=1$ and thus $\tau$ belongs to $(\FF^*_{\ell^2})^{\frac{d}{r}(\ell-1)}$, completing the proof.

\end{proof}

\end{document}
